\chapter{总结与展望}\label{Chap:chap5}

\section{本文工作总结}\label{Sec5.1}

离线强化学习通过从预先收集的静态数据集中学习策略，避免了在线交互的高昂成本与安全风险，在医疗诊断、自动驾驶、机器人控制、金融决策等实际应用场景中具有重要价值。然而，离线数据集的有限覆盖范围导致的分布偏移问题是制约该技术发展的核心挑战：当智能体的策略访问离线数据未覆盖的状态-动作对时，价值函数的估计将产生外推误差，导致策略性能显著低于预期。本文围绕离线强化学习中的不确定性估计与探索策略展开系统研究，针对纯离线设定和离线到在线设定分别提出了两种互补的算法方案，取得了以下研究成果：

第一，提出基于自适应不确定性度量的离线强化学习算法。针对现有基于模型方法仅考虑状态转移不确定性、采用静态惩罚系数的局限，本文从两个维度进行改进：在不确定性度量层面，引入累积状态不确定性的概念，通过递推公式$u(s') = \omega \cdot u(s) + \log(\Sigma_\theta(s, a) + 1)$综合考虑从初始状态到当前状态的历史不确定性累积效应，使智能体能够更全面地评估序列决策中的风险累积；在惩罚调节层面，设计基于Sigmoid函数的自适应动态权重机制$\lambda_k = \lambda_0 / (1 + e^{K_0 - k})$，根据策略优化进程动态调节不确定性惩罚强度，在训练初期降低惩罚以鼓励探索搜索策略空间，在策略逐渐收敛时增加惩罚以保持保守性确保策略稳定。理论分析证明了累积不确定性沿轨迹的单调不减性质和不确定性预测网络的泛化能力，以及自适应权重相比固定权重在探索-保守平衡方面的优势。在D4RL MuJoCo基准测试集上的实验表明，该方法在9个任务中的6个取得最优或接近最优性能，特别是在数据质量较差的随机数据集上相比MOPO取得了显著提升（halfcheetah任务从35.4提升至37.2，hopper任务从11.7提升至17.6），在hopper-medium任务上提升幅度达44.4个百分点。

第二，提出基于定向探索模型的离线到在线强化学习算法。针对现有离线到在线方法保守性与灵活性矛盾、不确定性利用不充分、在线交互成本高等问题，本文提出将不确定性信息从“惩罚信号”转变为“探索引导信号”的核心思想。在离线阶段，通过集成模型和不确定性惩罚进行保守策略优化，确保策略在模型预测可靠的区域内更新；在在线阶段，通过定向探索目标$\nabla_\varphi J_\pi(\varphi) = \mathbb{E}[\nabla_\varphi \log \pi_\varphi(a|s) \cdot (c \cdot u_{on}(s, a) + Q(s, a))]$显式引导策略探索高不确定性（转移模型覆盖之外）且高动作价值（接近任务目标）的区域。在技术层面，提出双重不确定性估计方法区分偶然不确定性和认知不确定性在不同阶段的作用：离线阶段使用集成模型的预测方差估计偶然不确定性以实现保守性，在线阶段同时使用预测方差和集成分歧度估计两类不确定性以引导探索。设计动态权重自适应机制$c = c_0 / (1 + e^{(k - k_0) \cdot \nu})$实现探索与利用的平滑过渡，在交互初期注重不确定性引导的探索，在交互后期转向价值引导的利用。理论分析证明了在分段常数函数假设下，定向探索策略相比均匀采样具有更优的收敛速率（指数从$-1/d$改善至$-1/(d-1)$）。实验结果表明，该方法仅需50k步在线交互即可在medium-replay数据集上显著优于所有基线方法，在所有9个任务上均优于MOPO，通过t-SNE可视化验证了定向探索确实聚焦于高不确定性、高价值区域。

两种方法形成递进关系：首先，本文提出基于自适应不确定性度量的离线强化学习算法，解决纯离线设定下的保守策略优化问题；其次，引入在线微调阶段的定向探索模型，将“被动惩罚分布外”转变为“主动探索分布外”。两种方法共同构成了从纯离线到离线在线联合优化的完整技术方案。


\section{进一步研究方向}\label{Sec5.2}

尽管本文提出的方法在实验中取得了良好效果，但仍存在一定局限性，有待进一步研究改进。在算法效率方面，本文方法需要额外训练不确定性预测网络，增加了模型复杂度和训练开销，未来可探索更轻量级的不确定性估计方法，如基于梯度信息的在线估计、基于核方法的非参数估计、或利用神经网络中间层特征的即时估计，在保持估计精度的同时降低计算成本；在理论基础方面，本文采用Sigmoid函数实现惩罚权重的平滑过渡，但最优的过渡曲线可能与任务特性相关，未来可研究基于元学习或贝叶斯优化的自动权重调节方法，从数据中自适应学习最优的探索-保守平衡策略。

在算法适用性方面，实验表明本文方法在medium-expert等高质量数据集上的提升有限，这是因为专家数据的高集中度限制了定向探索的空间，未来可引入行为克隆损失项或设计自适应的探索门控机制，使算法能够根据离线数据质量自动调节探索强度——在低质量数据上加强探索以突破次优策略的局限，在高质量数据上减少探索以避免干扰专家知识；在分布外检测方面，本文方法主要基于集成模型的不确定性估计，未来可探索将表征学习与不确定性估计相结合的方法，通过学习状态-动作空间的紧致表征，利用表征空间的几何结构（如到数据中心的距离）实现更精确的分布外检测。

在应用拓展方面，本文方法针对单一任务设计，在多任务学习或任务迁移场景中可能需要额外的适配机制，未来可研究跨任务的不确定性共享与迁移方法，使模型能够利用相关任务的经验加速学习；在实际应用中的安全约束方面，医疗、自动驾驶等高风险场景除了性能优化外还需满足安全约束，未来可将本文的不确定性估计方法与约束强化学习相结合，在保证安全边界的前提下实现策略优化，例如将不确定性估计融入安全约束的构建，在高不确定性区域施加更严格的安全裕度。综合各方面分析，离线强化学习作为连接历史数据与智能决策的重要桥梁，具有广阔的应用前景。本文的研究为不确定性驱动的离线策略优化提供了新的思路，特别是将不确定性从被动规避的风险转化为主动利用的探索信号，为该领域的后续研究开辟了新的方向。
